\chapter{Eosinophilia}

\section*{Definition}
Eosinophilia is an increased blood eosinophil count, usually defined as more than 500 cells per microliter. Persistent counts of 1,500 cells per microliter or higher, especially with organ symptoms, require evaluation for hypereosinophilic disease.

\section*{What Happens in Your Body}
Eosinophils participate in allergic inflammation and defense against some parasites. Excessive or persistent activation can injure the skin, lungs, gastrointestinal tract, heart, nervous system, or other organs.

\section*{Causes}
\begin{itemize}
  \item Atopic disease such as asthma, allergic rhinitis, eczema, and drug hypersensitivity.
  \item Parasitic infection, particularly with travel or exposure to contaminated food, soil, or water.
  \item Medication reactions, autoimmune or inflammatory disease, adrenal insufficiency, and selected cancers.
  \item Clonal or idiopathic hypereosinophilic syndromes.
\end{itemize}

\section*{When to See a Doctor}
Seek urgent care for chest pain, shortness of breath, neurologic symptoms, confusion, severe abdominal pain, rapidly spreading rash, facial swelling, or fever after starting a medicine. Persistent marked eosinophilia or evidence of heart, lung, or nerve injury needs prompt specialist assessment.

\section*{Related Symptoms}
\begin{itemize}
  \item Itching, rash, wheezing, sinus symptoms, cough, abdominal pain, diarrhea, fever, fatigue, or weight loss.
  \item Some people have no symptoms and are identified through a complete blood count.
\end{itemize}

\section*{Self-Care / Home Management}
Do not re-challenge yourself with a suspected medication. Record new medicines, supplements, foods, animal exposures, travel, and rash or respiratory symptoms. Avoid empiric antiparasitic treatment without appropriate clinical guidance.

\section*{How Your Doctor Will Evaluate You}
\begin{itemize}
  \item Repeat complete blood count with differential and examine the trend and absolute count.
  \item Review medicines, allergies, asthma, travel, exposures, and systemic symptoms; test stool or blood for parasites when appropriate.
  \item Persistent or marked eosinophilia may require imaging, organ-function testing, molecular studies, and hematology evaluation.
\end{itemize}

\section*{Treatment Options}
Treatment removes the trigger and manages the cause: stopping a culprit drug, targeted antiparasitic therapy, control of allergic disease, or corticosteroids and specialist-directed targeted therapy for organ-threatening hypereosinophilic disease.

\section*{References}
\begin{thebibliography}{99}
\bibitem{eosinophilia:ref1} Gotlib J. World Health Organization-defined eosinophilic disorders. \textit{Hematology Am Soc Hematol Educ Program}. 2017;2017:443--456.
\bibitem{eosinophilia:ref2} Rothenberg ME. Eosinophilia. \textit{N Engl J Med}. 1998;338:1592--1600.
\end{thebibliography}
